\documentclass[10pt,a4paper]{amsart}
\usepackage[margin=19mm]{geometry}
\usepackage{amsmath,amssymb,mathtools}
\usepackage[scr=rsfs,cal=euler]{mathalfa}
\usepackage{xfrac}
\usepackage[unicode,hidelinks,bookmarks=false]{hyperref}
\newcommand{\ie}{i.e.\ }
\newcommand{\cf}{cf.\ }
\newcommand{\resp}{respectively\ }
\newcommand{\Subsection}{\subsection}
\providecommand{\nobreakdash}{\nobreak\hbox{-}}
\setlength{\emergencystretch}{2em}
\multlinegap0pt
\usepackage{amssymb}
\newtheorem{lemma}{Lemma}
\title{Stochastic fractional heat equation with general rough noise}
\author{Bin Qian, Ran Wang}
\date{Source: arXiv:2604.07697v1 (CC BY 4.0)}
\begin{document}
\maketitle

\noindent\emph{Source and licence.} Stochastic fractional heat equation with general rough noise. Version arXiv:2604.07697v1 (CC BY 4.0), distributed by arXiv under the Creative Commons Attribution 4.0 International licence, http://creativecommons.org/licenses/by/4.0/, as recorded in the arXiv licence metadata for this exact version and shown on the abstract page https://arxiv.org/abs/2604.07697v1. Full licence notice: \texttt{licenses/arxiv-2604.07697-CC-BY-4.0.txt}.\par\smallskip

\noindent\textit{Editorial scope.} Counted unit: the article's lemma integ (source0.tex lines 553--557). The article's complete original proof of it is delivered with it (source0.tex lines 558--602, one \texttt{proof} environment of 45 lines). No mathematics is written, repaired or re-derived: the statement and the proof are the article's own lines, in the article's own notation, and no retained line is substituted or re-flowed.  No further source-local context is needed: every cross-reference of every retained line resolves inside the two delivered spans.  Nothing was omitted from any retained span, so the package carries no omission record.  No retained line contains a citation, so the wrapper carries no bibliography and no bibliography entry.  No retained line contains a figure environment, an image-inclusion command or drawing code, so no image is delivered and the package declares no asset dependency.  Editorial formatting only, in addition: an AMS \texttt{amsart} preamble that restates the article's own declarations and macros used by the retained lines, namely the environments \texttt{lemma} on separate counters, together with the standard packages those lines need.  The retained environments print in this document as Lemma 1 in order of appearance, whereas the article prints the same items with the numbering of its own full document.  The complete native source of the article as distributed in the arXiv e-print for this exact version is bundled unchanged as \texttt{sources/198108/source0.tex}.
\begin{lemma}\label{lem integ} For any $H\in (0,\frac12)$, there exists a constant  $c_H>0$ depending on $H$ such that for any $x\in\mathbb{R}$,
 \begin{equation}\label{eq integ aim0}
 \int_{|y|>1}|y|^{2H-2} \left(1\wedge |x+y|^{2H-2}\right)dy\le c_H\left(1\wedge |x|^{2H-2}\right).
\end{equation}
\end{lemma}

\begin{proof}
Since
 \begin{equation*}
 \int_{|y|>1}|y|^{2H-2} \left(1\wedge |x+y|^{2H-2}\right)dy\le \int_{|y|>1}|y|^{2H-2}dy=\frac{2}{1-2H},
\end{equation*} 
it suffices to show that for any $x>2$,
\begin{equation}\label{eq integ 1}
\int_{1}^{\infty}|y|^{2H-2}\left(1\wedge |x+y|^{2H-2}\right)dy\lesssim   |x|^{2H-2},
\end{equation}
\begin{equation}\label{eq integ 2}
\int_{1}^{\infty}|y|^{2H-2}\left(1\wedge |x-y|^{2H-2}\right)dy \lesssim   |x|^{2H-2}.
\end{equation}

Estimate \eqref{eq integ 1} follows easily from 
\begin{align*}
\int_{1}^{\infty}|y|^{2H-2}\left(1\wedge |x+y|^{2H-2}\right) dy \le \int_{1}^{\infty}|y|^{2H-2} x^{2H-2}dy\lesssim x^{2H-2}.
\end{align*}

It remains to prove \eqref{eq integ 2}  for any $x>2$. We decompose the integral as 
\begin{equation}\label{eq integ 3}
\begin{split}
&\int_{1}^{\infty}|y|^{2H-2}\left(1\wedge |x-y|^{2H-2}\right)dy\\
 =&\, \int_{1}^{\frac{x}2}y^{2H-2}\left(1\wedge |x-y|^{2H-2}\right)dy+\int_{\frac{x}{2}}^{\infty}y^{2H-2}\left(1\wedge |x-y|^{2H-2}\right)dy  \\
 =:&\, I_1+I_2.  
\end{split}
\end{equation}

For  $y\in [1,\frac{x}2]$,  we have $x-y\ge \frac{x}{2}>1$.  Hence, 
 \begin{align}\label{eq integ 31}
{I_1}\le4 x^{2H-2}\int_1^{\frac{x}2}y^{2H-2}dy\lesssim x^{2H-2}.
\end{align}

For    $y\in [\frac{x}2,\infty]$, we have   $y^{2H-2}\le 4x^{2H-2}$.  Hence, 
\begin{equation}\label{eq integ 32}
\begin{split}
{I_2} \le &\,  4 x^{2H-2}\int_{\frac{x}2}^{\infty}\left(1\wedge |x-y|^{2H-2}\right)dy\\
%&=4x^{2H-2}\int_{\{[\frac{x}2,x-1)\cup[x-1,x+1]\cup(x+1,\infty]\}}\left(1\wedge |x-y|^{2H-2}\right)dy\\
\le &\, 4x^{2H-2} \left(\int_{x-1}^{x+1}1dy+\int_{[\frac{x}2,x-1]\cup [x+1,\infty]}|x-y|^{2H-2}dy\right)\\
\le &\,  4x^{2H-2}\left(2+2\int_{z\ge1} z^{2H-2}dz\right)\\
 \lesssim &\, x^{2H-2},
 \end{split}
\end{equation}
Putting \eqref{eq integ 3}-\eqref{eq integ 32} together, we obtain \eqref{eq integ 2}.
 The proof is complete.  
  \end{proof}

\end{document}
